\documentclass[border=3pt]{standalone}
% 公共导言区 —— 所有 TikZ 图 \input 本文件后使用
\usepackage[UTF8, fontset=mac]{ctex}
\usepackage{tikz}
\usetikzlibrary{arrows.meta, positioning, calc, shapes.geometric, decorations.pathmorphing, backgrounds, fit}
\definecolor{zhusha}{HTML}{9B2D20}
\definecolor{mohei}{HTML}{1A1A1A}
\definecolor{shenhui}{HTML}{555555}
\definecolor{qianhui}{HTML}{F5F0EC}
\definecolor{lan}{HTML}{1F4E79}
\definecolor{qing}{HTML}{2E7D6E}
\tikzset{
  block/.style={draw=shenhui, fill=white, thick, rounded corners=1.5pt,
                font=\small, align=center, inner sep=4pt, minimum height=7mm},
  core/.style={block, fill=lan!12, draw=lan, font=\small\bfseries},
  periph/.style={block, fill=qing!10, draw=qing},
  power/.style={block, fill=zhusha!8, draw=zhusha},
  note/.style={font=\footnotesize, text=shenhui, align=center},
  lab/.style={font=\footnotesize, text=mohei},
  sig/.style={-{Stealth[length=2.2mm]}, thick, color=mohei},
  fbsig/.style={-{Stealth[length=2.2mm]}, thick, color=zhusha, dashed},
  vec/.style={-{Stealth[length=3mm]}, very thick, color=zhusha},
}

\begin{document}
\begin{tikzpicture}[x=1mm, y=1mm]

% ===== CPU 域（左）=====
\node[core, minimum width=34mm, minimum height=11mm] (m4) at (0,0) {Cortex-M4F\\170 MHz · FPU · DSP};
\node[block, below=3mm of m4, fill=zhusha!10, draw=zhusha] (cordic) {CORDIC\\三角函数加速};
\node[block, below=3mm of cordic, fill=zhusha!10, draw=zhusha] (fmac) {FMAC\\FIR/IIR 滤波加速};
\node[draw=lan, thick, rounded corners=2pt, fit=(m4)(cordic)(fmac), inner sep=3mm, fill=lan!4] (cpudom) {};
\node[note, anchor=south west] at (cpudom.north west) {数学加速器域};

% ===== 存储（中左）=====
\node[block] (flash) at (30mm, 8mm) {Flash\\128 KB 双 bank};
\node[block, below=3mm of flash] (sram) {SRAM 32 KB\\(含 CCM SRAM)};
\node[block, below=3mm of sram, fill=qianhui] (bus) {总线矩阵\\AHB / APB 桥};

% ===== 模拟外设（右）=====
\node[periph] (tim) at (60mm, 24mm) {TIM1/TIM8 高级定时器\\互补 PWM · 死区 · 刹车};
\node[periph, below=2.5mm of tim] (adc) {2× ADC 12-bit 4 MSPS\\双 ADC 同步 · 硬件过采样};
\node[periph, below=2.5mm of adc] (dac) {4 通道 DAC 12-bit};
\node[periph, below=2.5mm of dac] (opamp) {3× 运放（PGA 增益 2--16）};
\node[periph, below=2.5mm of opamp] (comp) {4× 超快比较器 → 刹车};
\node[periph, below=2.5mm of comp] (vref) {VREFBUF 基准};

% ===== 数字外设（右 2）=====
\node[block] (dma) at (94mm, 24mm) {DMA×2};
\node[block, below=2.5mm of dma] (fdcan) {FDCAN};
\node[block, below=2.5mm of fdcan] (usb) {USB FS};
\node[block, below=2.5mm of usb] (com) {SPI / I2C / USART};
\node[block, below=2.5mm of com] (gpio) {EXTI / GPIO};
\node[block, below=2.5mm of gpio, fill=qianhui] (dbg) {调试 SWD / ST-LINK};

% 连线
\draw[sig] (cpudom.east) -- (bus.west);
\draw[sig] (flash.south) -- (bus.north);
\draw[sig] (sram.north) -- ($(flash.south)+(6mm,0)$) -- ($(bus.north)+(6mm,0)$);
\draw[sig] (bus.east) -- ++(4mm,0) coordinate (bx);
\draw[sig, color=shenhui] (bx) -- (tim.west);
\draw[sig, color=shenhui] (bx) -- (adc.west);
\draw[sig, color=shenhui] (bx) -- (dac.west);
\draw[sig, color=shenhui] (bx) -- (opamp.west);
\draw[sig, color=shenhui] (bx) -- (comp.west);
\draw[sig, color=shenhui] (bx) -- (vref.west);
\draw[sig] (dma.west) -- ++(-4mm,0) coordinate (bx2);
\draw[sig, color=shenhui] (bx2) -- (fdcan.west);
\draw[sig, color=shenhui] (bx2) -- (usb.west);
\draw[sig, color=shenhui] (bx2) -- (com.west);
\draw[sig, color=shenhui] (bx2) -- (gpio.west);
\draw[sig, color=shenhui] (bx2) -- (dbg.west);

% 电机控制强调：TIM1 ↔ ADC ↔ COMP ↔ OPAMP 圈出
\node[draw=qing, thick, rounded corners=3pt, fit=(tim)(adc)(opamp)(comp), inner sep=2.5mm, dashed] (mc) {};
\node[note, text=qing, anchor=south] at (mc.north) {电机控制信号链（单芯片闭环）};

\end{tikzpicture}
\end{document}
